
dans les deux cas envisag\'es. Comme @@M0@@ est propre, @@M1@@ \'etant normal (FGA 236-2.13), et comme @@M2@@ est un sous-groupe, n\'ecessairement ferm\'e (SGA 3 VI@@M3@@ 1.9.2), pour que @@M4@@ soit fini, il suffit qu'il ne contienne pas le sous-groupe de type @@M5@@ o\`u @@M6@@.

Consid\'erons la th\'eorie de Kummer, c'est-\`a-dire, le diagramme commutatif suivant, provenant de la suite exacte
@@M7@@
de faisceaux \'etales sur @@M8@@ (resp. sur @@M9@@)~:
@@M10@@
Par le ``th\'eor\`eme 90'' de Hilbert, on a @@M11@@ (resp. @@M12@@). Comme @@M13@@ est propre et int\`egre, on a @@M14@@\,; donc @@M15@@ est injective.

Il suffit donc de voir que @@M16@@ est injectif, ce qui r\'esulte du fait que @@M17@@ est surjectif, ou encore (SGA 1 V 6.9) que tout rev\^etement \'etale connexe @@M18@@ de @@M19@@ de groupe @@M20@@ a une restriction @@M21@@ qui est aussi connexe, ce qui est une cons\'equence bien connue du th\'eor\`eme de Bertini et du th\'eor\`eme de connexion de Zariski (SGA 1 X 2.10).

\textbf{Remarque 3.12.} Conservons les conditions de 3.11. Alors :
\begin{enumerate}[label=(\roman*)]
\item Le noyau @@M22@@ est un groupe fini \uline{unipotent} (SGA 3 XVII 1.1), puisqu'il ne contient pas le sous-groupe de type multiplicatif @@M23@@ pour aucun @@M24@@ (SGA 3 XVII 4.6.1 vi), ce qu'on prouve de m\^eme, utilisant au
